\section{Sliding Window Maximum}
\label{sec:sq-sliding-window-max}

\problemstatement{Given an array $\textit{nums}$ and a window size $k$,
return an array containing the maximum element of every contiguous
window of size $k$ as it slides from the start of the array to the end.}

\begin{patternbox}
\emph{``Maximum (or minimum) over every sliding window''} is the
\textbf{monotonic deque} keyword from Section~\ref{sec:sq-intro}: a
single index can remain a useful candidate for the maximum across
\emph{many} consecutive windows, but it must also eventually be evicted
once the window slides past it -- requiring removal from \textbf{both}
ends (the back, when a better candidate arrives; the front, when the
oldest index falls out of the window), which is exactly what a
double-ended queue (deque) provides.
\end{patternbox}

\subsection*{Understanding the problem carefully}

Maintain a deque of \emph{indices} (not values), kept in decreasing order
of their corresponding values from front to back, and restricted to
indices currently inside the window. The front of the deque is always the
index of the current window's maximum, because it is both the largest
remaining value (by the decreasing-order invariant) and guaranteed to
still be within the window (indices that fall out are actively removed
from the front).

\begin{ideabox}[Monotonic Deque: Evict from the Back on Value, from the Front on Position]
For each new index $i$: first, while the deque is non-empty and the value
at its back index is $\le \textit{nums}[i]$, pop from the back (those
indices can never again be the maximum of any window containing $i$,
since $i$ is both newer and at least as large -- the same domination
argument as next greater element). Push $i$ to the back. Then, while the
deque's front index is $\le i-k$ (i.e.\ has fallen outside the current
window), pop from the front. Once $i \ge k-1$, the deque's front index is
the maximum of the current window.
\end{ideabox}

\subsection*{Why eviction from both ends preserves correctness}

Eviction from the \textbf{back} is exactly the monotonic-stack domination
argument from Section~\ref{sec:sq-next-greater-1}: a smaller, older index
can never be the maximum once a newer, at-least-as-large value has
arrived, anywhere within the window. Eviction from the \textbf{front} is
new to this section and is required because, unlike a plain monotonic
stack, indices here have a finite lifetime determined by the window's
position -- even an index that has never been dominated by value must
still be discarded once it ages out of the window. Both evictions
together guarantee that the deque, at every step, contains exactly the
indices that are both currently in-window \emph{and} not dominated by
any later in-window value, kept in decreasing order -- so its front is
always correct.

\subsection*{Algorithm}

\begin{enumerate}
  \item Initialize an empty deque of indices and an empty result list.
  \item For $i$ from $0$ to $n-1$:
  \begin{enumerate}
    \item While the deque is non-empty and
          $\textit{nums}[\text{back index}] \le \textit{nums}[i]$: pop
          from the back.
    \item Push $i$ to the back.
    \item While the deque's front index $\le i-k$: pop from the front.
    \item If $i \ge k-1$: append $\textit{nums}[\text{front index}]$ to
          the result.
  \end{enumerate}
  \item Return the result.
\end{enumerate}

\subsection*{Dry run}

\dryrun{Take $\textit{nums} = [1,3,-1,-3,5,3,6,7]$, $k=3$.}

\begin{itemize}
  \item $i=0$ ($1$): deque empty, push. Deque $=[0]$. $i<k-1=2$, no
        output.
  \item $i=1$ ($3$): back is index $0$ (value $1\le3$), pop. Push $1$.
        Deque $=[1]$. $i<2$, no output.
  \item $i=2$ ($-1$): back is index $1$ (value $3 \not\le -1$), no pop.
        Push $2$. Deque $=[1,2]$. Front index $1 > i-k=-1$, no pop.
        $i=2\ge2$: output $\textit{nums}[1]=3$.
  \item $i=3$ ($-3$): no pop (back value $-1\not\le-3$). Push $3$.
        Deque $=[1,2,3]$. Front index $1 \le i-k=0$? No ($1>0$). Output
        $\textit{nums}[1]=3$.
  \item $i=4$ ($5$): back values $-3,-1,3$ all $\le5$: pop all. Deque
        $=[]$, push $4$. Deque $=[4]$. Front $4 \le i-k=1$? No. Output
        $\textit{nums}[4]=5$.
  \item $i=5$ ($3$): back value $5\not\le3$, no pop. Push $5$. Deque
        $=[4,5]$. Front $4\le i-k=2$? No. Output $\textit{nums}[4]=5$.
  \item $i=6$ ($6$): back values $3,5$ both $\le6$: pop both. Push $6$.
        Deque $=[6]$. Front $6\le i-k=3$? No. Output
        $\textit{nums}[6]=6$.
  \item $i=7$ ($7$): back value $6\le7$, pop. Push $7$. Deque $=[7]$.
        Front $7\le i-k=4$? No. Output $\textit{nums}[7]=7$.
\end{itemize}

Result: $[3,3,5,5,6,7]$.

\subsection*{C++ implementation}

\begin{lstlisting}
#include <bits/stdc++.h>
using namespace std;

vector<int> maxSlidingWindow(vector<int>& nums, int k) {
    deque<int> dq; // stores indices
    vector<int> result;

    for (int i = 0; i < (int)nums.size(); i++) {
        while (!dq.empty() && nums[dq.back()] <= nums[i]) {
            dq.pop_back();
        }
        dq.push_back(i);

        while (dq.front() <= i - k) {
            dq.pop_front();
        }

        if (i >= k - 1) {
            result.push_back(nums[dq.front()]);
        }
    }
    return result;
}

int main() {
    vector<int> nums = {1, 3, -1, -3, 5, 3, 6, 7};
    for (int x : maxSlidingWindow(nums, 3)) cout << x << " ";
    cout << endl;
    return 0;
}
\end{lstlisting}

\begin{complexitybox}
\textbf{Time Complexity.} Each index is pushed to the back once and
popped (from either end) at most once across the whole scan, giving
\[
  O(n).
\]
\textbf{Space Complexity.} $O(k)$ for the deque in the worst case.
\end{complexitybox}

\begin{takeawaybox}
A monotonic deque is a monotonic stack with one extra capability: aging
out elements from the opposite end once they fall outside a window. Any
time a problem combines ``sliding window'' with ``maximum/minimum within
the window,'' this is the data structure to reach for, rather than
re-scanning each window from scratch (which would cost $O(nk)$) or using
a heap (which would cost $O(n \log k)$ but require extra bookkeeping to
evict stale entries lazily).
\end{takeawaybox}
